\documentclass[11pt]{article}
\usepackage{amsmath,amssymb,amsthm}
\usepackage[margin=1.2in]{geometry}
\usepackage{hyperref}

\newtheorem{theorem}{Theorem}
\newtheorem{lemma}[theorem]{Lemma}
\newtheorem{corollary}[theorem]{Corollary}
\newtheorem{conjecture}[theorem]{Conjecture}
\theoremstyle{remark}
\newtheorem{remark}[theorem]{Remark}

\newcommand{\Zi}{\mathbb{Z}[i]}
\newcommand{\Z}{\mathbb{Z}}
\newcommand{\N}{\mathbb{N}}
\newcommand{\nup}{\nu_p}
\DeclareMathOperator{\im}{Im}
\DeclareMathOperator{\re}{Re}

\title{A valuation obstruction for $3\times 3$ magic squares of squares,\\
with a machine-checked proof}
\author{Florian Ernst \and Albin Morel}
\date{August 2026 --- draft}

\begin{document}
\maketitle

\begin{abstract}
We study the open problem of whether a $3\times 3$ magic square whose nine
distinct entries are perfect squares exists.  We reduce the problem to a purely
additive condition on the set $D(e)=\{2xy : x^2+y^2=e^2\}$ attached to the
center root $e$, and prove unconditionally that the center root of any such
square must be divisible by at least two distinct primes $p\equiv 1\pmod 4$:
for centers $e = s\,p^a$ with a single such prime, the elements of $D(e)$ have
pairwise distinct $p$-adic valuations, which is incompatible with the required
additive structure.  All results, including the reduction, are formally
verified in Lean~4 with mathlib.  We also report an exhaustive computation
showing no magic square of squares has center entry below $10^{20}$, and that
Bremner's fully magic square with seven square entries is, up to scaling, the
only one with center root below $10^9$.  Finally we describe a rigidity
phenomenon that pins down the remaining two-prime case to explicit divisibility
coincidences.
\end{abstract}

\section{Introduction}

Whether nine distinct perfect squares can form a fully magic $3\times 3$
square is a well-known open problem, popularized by Martin Gardner's 1996
prize and studied extensively since; see Boyer's survey and problem
collection~\cite{boyer}.  The strongest results to date are computational
(searches over magic sums and over arithmetic-progression parameters) together
with structural constraints on hypothetical solutions due to Morgenstern,
Roberts, Underwood and others~\cite{boyer}.

Throughout, a \emph{magic square of squares} is a $3\times 3$ array of nine
distinct perfect squares whose three rows, three columns and both diagonals
share a common sum $S$.  Summing the four lines through the center shows
$S = 3E$ where $E=e^2$ is the center entry; we call $e$ the \emph{center
root}.

\subsection*{Results}

\begin{theorem}[Reduction]\label{thm:reduction}
Let $D(e)=\{2xy : x^2+y^2=e^2,\ 0<x<y\}$.  A magic square of squares with
center $e^2$ exists if and only if there are $u,v$ with
$u,\,v,\,u+v,\,u-v \in D(e)$ (all nonzero, $u\neq v$); the square's entries
are then $e^2,\ e^2\pm u,\ e^2\pm v,\ e^2\pm(u+v),\ e^2\pm(u-v)$.
\end{theorem}

\begin{theorem}[Main theorem]\label{thm:main}
If the center root $e$ is divisible by at most one prime $p\equiv 1\pmod 4$
(to any power, and with arbitrary cofactor free of such primes), then no magic
square of squares with center $e^2$ exists.  Equivalently: the center root of
any magic square of squares is divisible by at least two distinct primes
$\equiv 1 \pmod 4$.
\end{theorem}

\begin{theorem}[Computations]\label{thm:comp}
There is no magic square of squares with center entry $e^2 < 10^{20}$.
Moreover, every fully magic $3\times 3$ square with exactly seven square
entries and center root $e<10^9$ is an integer scaling of Bremner's example
with $e=425$.
\end{theorem}

All statements of Theorems~\ref{thm:reduction} and~\ref{thm:main}, including
Fermat's theorem that three squares in arithmetic progression have common
difference divisible by $24$, are formally verified in Lean~4 over
mathlib; the library (thirteen files, no unproven assumptions) accompanies
this paper.

\section{The reduction}

Every $3\times 3$ magic square has the form
$c \pm \{0,u,v,u+v,u-v\}$ arranged as
\[
\begin{pmatrix} c+u & c-u-v & c+v\\ c-u+v & c & c+u-v\\ c-v & c+u+v & c-u
\end{pmatrix},
\]
in which all eight line sums equal $3c$ identically.  Thus magic is free and
the entire problem is the squareness of the nine cells.

\begin{lemma}\label{lem:pair}
For $c=e^2$, both $c-d$ and $c+d$ are squares if and only if $d\in D(e)\cup\{0\}$
(up to sign).  Indeed if $x^2+y^2=e^2$ then $c-2xy=(x-y)^2$ and
$c+2xy=(x+y)^2$; conversely if $c-d=A^2$ and $c+d=B^2$ then $A^2+B^2=2e^2$,
$A\equiv B \pmod 2$, and $x=\tfrac{B-A}2$, $y=\tfrac{B+A}2$ satisfy
$x^2+y^2=e^2$, $2xy=d$.
\end{lemma}

Theorem~\ref{thm:reduction} follows by applying Lemma~\ref{lem:pair} to the
four center lines.  (Distinctness of the nine entries corresponds to
$u,v,u\pm v \neq 0$ and $u \neq v$.)

\section{Proof of the main theorem}

Fix a prime $p\equiv 1\pmod 4$, write $p=\pi\bar\pi$ in $\Zi$ with
$\pi = A + Bi$, $A^2+B^2=p$, and let $e = s\,p^a$ with $a\ge 0$ and $s$
free of primes $\equiv 1\pmod 4$.

\begin{lemma}[Rigid part]\label{lem:rigid}
Every $z\in\Zi$ with $N(z)=s^2$ is a unit multiple of $s$.
\end{lemma}

\begin{proof}
Induct on the least prime factor $q$ of $s$.  If $q\equiv 3\pmod 4$ then $q$
is inert; $q \mid N(z) = z\bar z$ gives $q\mid z$, and we divide.  If $q=2$,
then $1+i$ divides $z$ twice (as $2\mid N(z)$ twice), and
$(1+i)^2 = 2i$ is $2$ times a unit.
\end{proof}

\begin{lemma}[Classification]\label{lem:classify}
Every $z \in \Zi$ with $N(z) = s^2 p^{\,n}$ equals
$u\, s\, \pi^{j} \bar\pi^{\,n-j}$ for a unit $u$ and some $0\le j\le n$.
\end{lemma}

\begin{proof}
Induct on $n$: for $n\ge 1$, $\pi \mid p \mid z\bar z$, so $\pi\mid z$ or
$\bar\pi \mid z$; divide and recurse, finishing with Lemma~\ref{lem:rigid}.
\end{proof}

\begin{lemma}[Bridge]\label{lem:bridge}
For every $m\ge 1$, $p \nmid \im(\pi^m)$.
\end{lemma}

\begin{proof}
Pick $\omega \in \Z/p$ with $\omega^2=-1$ (possible as $p\equiv 1 \bmod 4$)
and consider the ring maps $\varphi_{\pm}:\Zi\to\Z/p$, $i\mapsto\pm\omega$.
If $p \mid \im(\pi^m)$ then, since
$\re(\pi^m)^2+\im(\pi^m)^2=p^m$, also $p\mid\re(\pi^m)$, so
$\varphi_\pm(\pi)^m=0$; as $\Z/p$ is a field, $\varphi_\pm(\pi)=0$, i.e.
$A\pm B\omega \equiv 0$.  Adding and subtracting gives $p\mid A$, $p \mid B$,
contradicting $A^2+B^2=p$.
\end{proof}

\begin{lemma}[Representation structure]\label{lem:repstruct}
Let $x^2+y^2 = e^2 = s^2p^{2a}$ with $xy\neq 0$.  Then there are
$1\le t\le a$ and $\varepsilon\in\{\pm1\}$ with
\[
2xy \;=\; \varepsilon\; p^{\,2(a-t)}\;\bigl(s^2 \im \pi^{4t}\bigr),
\qquad p \nmid s^2 \im \pi^{4t} .
\]
In particular $\nup(2xy) = 2(a-t)$, and distinct elements of $D(e)$ have
pairwise distinct $p$-adic valuations $\{0,2,\dots,2a-2\}$.
\end{lemma}

\begin{proof}
Write $z=x+yi$; by Lemma~\ref{lem:classify}, $z=u s\pi^j\bar\pi^{2a-j}$.
Then $2xy=\im(z^2)$ and $z^2=u^2 s^2 \pi^{2j}\bar\pi^{2(2a-j)}$ with
$u^2=\pm 1$.  If $j=a$ the product is real and $xy=0$.  Otherwise, with
$t=|j-a|$, $\pi^{2j}\bar\pi^{2(2a-j)} = p^{2(a-t)}\cdot\pi^{\pm4t}$, and
$\im(\bar\pi^{4t})=-\im(\pi^{4t})$; absorb signs into $\varepsilon$.
Lemma~\ref{lem:bridge} and $p\nmid s$ give the non-divisibility.
\end{proof}

\begin{proof}[Proof of Theorem~\ref{thm:main}]
Suppose $u,v,u+v,u-v\in D(e)$ as in Theorem~\ref{thm:reduction}, and give
each the form of Lemma~\ref{lem:repstruct} with parameters
$(t_i,\varepsilon_i)$ and constants $c_t = s^2\im\pi^{4t}$.
If $t_1=t_2$ then $u=\pm v$, excluded.  Otherwise
$\nup(u)\neq\nup(v)$; writing $u\pm v = p^{\,\min}\cdot(\text{bracket})$ where
the bracket is $\varepsilon_i c_{t_i} + (\text{multiple of }p)$, the bracket
is prime to $p$.  Comparing with the exact forms of $u+v$ and $u-v$ and using
uniqueness of exact $p$-power extraction, $2(a-t_3)=2(a-t_4)=\min$, so
$t_3=t_4$ and $u+v=\pm(u-v)$, forcing $u=0$ or $v=0$.  Contradiction.
\end{proof}

\begin{remark}
The counting argument (four center lines require four representations)
already forces two distinct primes \emph{or} $p^4\mid e$; the valuation
argument closes the prime-power loophole entirely and is what we formalize.
\end{remark}

\section{Computations}

Enumerating candidate centers directly through
Theorem~\ref{thm:reduction} is far cheaper than enumerating magic sums:
factoring $e$ in windows and generating $D(e)$ from the Gaussian
factorization checks all $e<10^{10}$ (center entries to $10^{20}$) in minutes
on a desktop; no configuration $u,v,u\pm v \in D(e)$ exists in that range.
The same machinery classifies near-misses: configurations with two of the
four differences in $D(e)$ and the other two ``half in'' (one square cell
each) are exactly the fully magic squares with seven square entries, of which
Bremner's $e=425$ example was the only one known.  Scanning all $e<10^9$
finds $2{,}352{,}941$ such configurations --- every single one an integer
scaling of $e=425$, and none with eight square entries.

\section{The two-prime frontier}

For $e = s\,p^aq^b$ with two distinct primes $\equiv 1\pmod 4$, the
valuation pairs $(\nup,\nu_q)$ organize $D(e)$ into \emph{conjugate pairs}
$\{d_{j,k}, d_{j,-k}\}$ of multiplicity exactly two off the axes, and axis
singletons.  If $\{u+v,u-v\}$ is a conjugate pair $\{d,d'\}$ arising from
$A=\pi^{4j}$, $B=\chi^{4k}$, then
\[
d+d' = 2\,\mathrm{scale}\cdot|\re B \,\im A|,\qquad
d-d' = 2\,\mathrm{scale}\cdot|\im B\, \re A|,
\]
so $u$ and $v$ are \emph{completely determined} --- and their valuation pair
coincides with that of $\{d,d'\}$, forcing $u,v\in\{d,d'\}$, which is absurd.
Consequently every hypothetical two-prime solution requires an explicit
divisibility coincidence ($q^t \mid \re \pi^{4j}$, $p^t\mid\re\chi^{4k}$, or a
mod-$p$ or mod-$q$ cancellation in $u\pm v$).  These coincidences recur along
arithmetic progressions of exponents, so closing them requires
exact-magnitude input beyond valuations; we leave this as the principal open
direction, noting that together with Theorem~\ref{thm:main} it would force
every solution's center root into three or more primes $\equiv 1\pmod 4$.

\begin{thebibliography}{9}
\bibitem{boyer} C.~Boyer, \emph{Latest research on the ``3x3 magic square of
squares'' problem}, \url{http://www.multimagie.com/English/SquaresOfSquaresSearch.htm}.
\bibitem{bremner} A.~Bremner, \emph{On squares of squares~II}, Acta Arith.
99 (2001), 289--308.
\bibitem{gardner} M.~Gardner, \emph{The magic of 3x3}, Quantum 6 (1996), 24--26.
\end{thebibliography}

\end{document}
